%======================================================
% CHAPTER V: SWT - SOFTWARE TESTING DOCUMENTATION
%======================================================
\chapter{SWT -- TÀI LIỆU KIỂM THỬ PHẦN MỀM}

\section{Tổng quan Tài liệu Kiểm thử}

Chương này trình bày \textbf{Tài liệu Kiểm thử Phần mềm (Software Testing Documentation -- SWT)} cho hệ thống Chatbot RAG. Mục tiêu của kiểm thử không chỉ đảm bảo các chức năng hoạt động đúng theo SWR, mà quan trọng hơn là đánh giá định lượng chất lượng của hệ thống AI thông qua bộ tiêu chuẩn RAGAS.

\begin{table}[H]
  \centering
  \caption{Thông tin tài liệu SWT}
  \label{tab:swt-info}
  \renewcommand{\arraystretch}{1.4}
  \begin{tabularx}{\textwidth}{|l|X|}
    \hline
    \rowcolor{blue!10}
    \textbf{Mục} & \textbf{Nội dung} \\
    \hline
    Phiên bản & 1.0 \\
    \hline
    Ngày kiểm thử & 06--07/2026 \\
    \hline
    Người kiểm thử & Trần Thị Kiều Trinh (Lead), hỗ trợ bởi các thành viên nhóm \\
    \hline
    Công cụ kiểm thử & Pytest, Postman, RAGAS Framework, Selenium (Frontend) \\
    \hline
  \end{tabularx}
\end{table}

%------------------------------------------------------
\section{Phạm vi Kiểm thử (Testing Scope)}
%------------------------------------------------------

\subsection{Trong phạm vi (In-scope)}

\begin{itemize}[itemsep=3pt]
  \item Kiểm thử chức năng (Functional Testing) cho tất cả Use Case trong SWR: UC-01 đến UC-10.
  \item Kiểm thử API (Backend) bằng Postman cho tất cả endpoint FastAPI.
  \item Kiểm thử tích hợp (Integration Testing) giữa Backend -- ChromaDB -- LLM.
  \item Kiểm thử hiệu năng (Performance Testing) đo Latency của RAG Pipeline.
  \item Kiểm thử AI (RAGAS Benchmark) đánh giá Accuracy, Faithfulness, Relevancy.
  \item Kiểm thử phi chức năng: xử lý lỗi, thông báo hệ thống, phân quyền.
\end{itemize}

\subsection{Ngoài phạm vi (Out-of-scope)}

\begin{itemize}[itemsep=3pt]
  \item Kiểm thử bảo mật chuyên sâu (Penetration Testing).
  \item Kiểm thử tải (Load Testing) quy mô lớn ($> 100$ người dùng đồng thời).
  \item Kiểm thử trên thiết bị di động.
  \item Kiểm thử khả năng phục hồi sau thảm họa (Disaster Recovery).
\end{itemize}

%------------------------------------------------------
\section{Chiến lược Kiểm thử (Testing Strategy)}
%------------------------------------------------------

\subsection{Mức độ Kiểm thử (Testing Levels)}

\begin{table}[H]
  \centering
  \caption{Các mức độ kiểm thử áp dụng}
  \label{tab:test-levels}
  \renewcommand{\arraystretch}{1.5}
  \begin{tabularx}{\textwidth}{|l|X|l|}
    \hline
    \rowcolor{blue!10}
    \textbf{Mức độ} & \textbf{Mô tả} & \textbf{Công cụ} \\
    \hline
    Unit Testing & Kiểm thử từng hàm/module độc lập (RAG Service, Embedding, Chunking). & Pytest \\
    \hline
    Integration Testing & Kiểm thử luồng dữ liệu giữa các module (Backend → ChromaDB → LLM). & Pytest + Postman \\
    \hline
    System Testing & Kiểm thử toàn bộ hệ thống theo kịch bản người dùng. & Postman, Manual \\
    \hline
    AI Quality Testing & Đánh giá chất lượng câu trả lời AI bằng RAGAS Benchmark. & RAGAS Framework \\
    \hline
    UAT & Người dùng thực tế (sinh viên) dùng thử và phản hồi. & Manual \\
    \hline
  \end{tabularx}
\end{table}

\subsection{Loại Kiểm thử (Testing Types)}

\begin{itemize}[itemsep=4pt]
  \item \textbf{Kiểm thử chức năng (Functional):} Xác nhận từng chức năng hoạt động đúng theo đặc tả SWR.
  \item \textbf{Kiểm thử hiệu năng (Performance):} Đo Latency của RAG Pipeline dưới điều kiện bình thường.
  \item \textbf{Kiểm thử phân quyền (Authorization):} Xác nhận Sinh viên không thể thực hiện tác vụ Admin.
  \item \textbf{Kiểm thử xử lý lỗi (Error Handling):} Xác nhận hệ thống trả về thông báo lỗi đúng chuẩn.
  \item \textbf{Kiểm thử biên (Boundary):} Kiểm thử file upload vượt 50MB, Top-K = 0, câu hỏi rỗng.
\end{itemize}

%------------------------------------------------------
\section{Kế hoạch Kiểm thử (Test Plan)}
%------------------------------------------------------

\subsection{Bảng Test Cases -- Kiểm thử Chức năng}

\begin{longtable}{|c|p{3.5cm}|p{3cm}|p{3cm}|c|}
  \caption{Test Cases kiểm thử chức năng hệ thống} \label{tab:test-cases} \\
  \hline
  \rowcolor{blue!10}
  \textbf{TC-ID} & \textbf{Tên Test Case} & \textbf{Input} & \textbf{Expected Output} & \textbf{Kết quả} \\
  \hline
  \endfirsthead
  \hline
  \rowcolor{blue!10}
  \textbf{TC-ID} & \textbf{Tên Test Case} & \textbf{Input} & \textbf{Expected Output} & \textbf{Kết quả} \\
  \hline
  \endhead
  \hline
  \endfoot
  %--- Chat RAG ---
  TC-01 & Chat RAG: Câu hỏi hợp lệ & Câu hỏi có trong tài liệu & Câu trả lời chính xác + nguồn trích dẫn & \textcolor{green!60!black}{\textbf{Pass}} \\
  \hline
  TC-02 & Chat RAG: Câu hỏi ngoài phạm vi & Câu hỏi không liên quan tài liệu & Thông báo MSG-03 & \textcolor{green!60!black}{\textbf{Pass}} \\
  \hline
  TC-03 & Chat RAG: Câu hỏi rỗng & Chuỗi rỗng & Thông báo lỗi validation & \textcolor{green!60!black}{\textbf{Pass}} \\
  \hline
  TC-04 & Chat RAG: Duy trì ngữ cảnh & Câu hỏi liên quan câu trước & Trả lời có ngữ cảnh từ lịch sử & \textcolor{green!60!black}{\textbf{Pass}} \\
  \hline
  TC-05 & Chat RAG: Thay đổi Top-K & Top-K = 1 / 5 / 10 & Số lượng sources trả về đúng & \textcolor{green!60!black}{\textbf{Pass}} \\
  \hline
  %--- Upload ---
  TC-06 & Upload: File PDF hợp lệ & File PDF $\leq$ 50MB & Status = Processing → Ready & \textcolor{green!60!black}{\textbf{Pass}} \\
  \hline
  TC-07 & Upload: File DOCX hợp lệ & File DOCX $\leq$ 50MB & Status = Processing → Ready & \textcolor{green!60!black}{\textbf{Pass}} \\
  \hline
  TC-08 & Upload: File sai định dạng & File .jpg, .xlsx & Thông báo MSG-02 & \textcolor{green!60!black}{\textbf{Pass}} \\
  \hline
  TC-09 & Upload: File quá lớn & File $>$ 50MB & Thông báo MSG-07 & \textcolor{green!60!black}{\textbf{Pass}} \\
  \hline
  TC-10 & Upload: File PDF bị hỏng & File PDF corrupt & Status = Error + thông báo & \textcolor{green!60!black}{\textbf{Pass}} \\
  \hline
  %--- Quản lý tài liệu ---
  TC-11 & Xóa tài liệu & Chọn tài liệu đã upload & Tài liệu bị xóa khỏi cả MySQL và ChromaDB & \textcolor{green!60!black}{\textbf{Pass}} \\
  \hline
  TC-12 & Xem danh sách tài liệu & Không có filter & Hiển thị đúng danh sách + metadata & \textcolor{green!60!black}{\textbf{Pass}} \\
  \hline
  %--- Phân quyền ---
  TC-13 & Phân quyền: Sinh viên xóa tài liệu & Sinh viên gọi DELETE /documents/1 & HTTP 403 Forbidden & \textcolor{green!60!black}{\textbf{Pass}} \\
  \hline
  TC-14 & Phân quyền: Sinh viên chạy Benchmark & Sinh viên gọi POST /benchmark/run & HTTP 403 Forbidden & \textcolor{green!60!black}{\textbf{Pass}} \\
  \hline
  %--- Benchmark ---
  TC-15 & Benchmark: Chạy kịch bản hoàn chỉnh & 50 câu hỏi, Recursive + Gemini Embedding & Điểm RAGAS được tính và lưu vào DB & \textcolor{green!60!black}{\textbf{Pass}} \\
  \hline
  TC-16 & Benchmark: Dashboard cập nhật & Sau khi benchmark xong & Dashboard hiển thị biểu đồ Radar mới & \textcolor{green!60!black}{\textbf{Pass}} \\
  \hline
  %--- Hiệu năng ---
  TC-17 & Hiệu năng: Latency câu hỏi thông thường & 1 câu hỏi, mạng ổn định & Latency $< 5$ giây & \textcolor{green!60!black}{\textbf{Pass}} \\
  \hline
  TC-18 & Hiệu năng: Ingestion tài liệu 50 trang & File PDF 50 trang & Ingestion hoàn tất $< 2$ phút & \textcolor{green!60!black}{\textbf{Pass}} \\
  \hline
  %--- Lỗi ---
  TC-19 & Xử lý lỗi: LLM API không phản hồi & Mock Gemini timeout & Thông báo MSG-04, không crash & \textcolor{green!60!black}{\textbf{Pass}} \\
  \hline
  TC-20 & Xử lý lỗi: ChromaDB không kết nối & Mock ChromaDB down & Thông báo lỗi kết nối DB & \textcolor{green!60!black}{\textbf{Pass}} \\
  \hline
\end{longtable}

%------------------------------------------------------
\section{Môi trường Kiểm thử Benchmark (RAGAS)}
%------------------------------------------------------

\subsection{Cấu hình Môi trường}

\begin{table}[H]
  \centering
  \caption{Cấu hình môi trường kiểm thử Benchmark}
  \label{tab:bench-env}
  \renewcommand{\arraystretch}{1.4}
  \begin{tabularx}{\textwidth}{|l|X|}
    \hline
    \rowcolor{blue!10}
    \textbf{Thành phần} & \textbf{Cấu hình} \\
    \hline
    Phần cứng & CPU: Intel Core i7-12700H, RAM: 16GB, GPU: NVIDIA RTX 3060 \\
    \hline
    LLM (cố định) & Gemini 2.5 Flash (via Google AI API) \\
    \hline
    Chunking Strategies & Recursive Character Text Splitter \textit{vs} Fixed-size Splitter \\
    \hline
    Chunk Size & 1000 ký tự, Overlap: 200 ký tự (đồng nhất cho cả hai) \\
    \hline
    Embedding Models & Google Gemini Embedding \textit{vs} BAAI/bge-m3 (local) \\
    \hline
    Top-K & 5 (mặc định cho toàn bộ kịch bản) \\
    \hline
    Test Dataset & 50 câu hỏi + Ground Truth trích từ giáo trình Công nghệ Phần mềm \\
    \hline
    Framework đánh giá & RAGAS (Python library) \\
    \hline
  \end{tabularx}
\end{table}

\subsection{Tập Dữ liệu Kiểm thử (Test Dataset)}

Bộ test dataset gồm 50 câu hỏi đa dạng, được phân loại theo độ khó:

\begin{table}[H]
  \centering
  \caption{Phân loại câu hỏi trong Test Dataset}
  \label{tab:dataset}
  \renewcommand{\arraystretch}{1.4}
  \begin{tabularx}{\textwidth}{|l|c|X|}
    \hline
    \rowcolor{blue!10}
    \textbf{Loại câu hỏi} & \textbf{Số lượng} & \textbf{Ví dụ} \\
    \hline
    Nhận biết định nghĩa & 15 & ``RAG là gì?'' \\
    \hline
    Hiểu biết quy trình & 15 & ``Quy trình Ingestion Pipeline gồm những bước nào?'' \\
    \hline
    Suy luận \& ứng dụng & 12 & ``Tại sao Recursive Chunking cho kết quả tốt hơn Fixed-size?'' \\
    \hline
    Tổng hợp nhiều nguồn & 8 & ``So sánh ưu nhược điểm của Gemini Embedding và BAAI/bge-m3.'' \\
    \hline
    \textbf{Tổng} & \textbf{50} & \\
    \hline
  \end{tabularx}
\end{table}

%------------------------------------------------------
\section{Kết quả Kiểm thử (Test Reports)}
%------------------------------------------------------

\subsection{Kết quả RAGAS Benchmark}

\begin{figure}[H]
  \centering
  \includegraphics[width=\textwidth]{images/test_benchmark.jpg}
  \caption{Giao diện cấu hình và khởi chạy Benchmark}
  \label{fig:benchmark-config}
\end{figure}

\begin{figure}[H]
  \centering
  \includegraphics[width=\textwidth]{images/dashboard.jpg}
  \caption{Dashboard tổng quan kết quả RAGAS sau Benchmark}
  \label{fig:benchmark-result}
\end{figure}

\begin{figure}[H]
  \centering
  \includegraphics[width=\textwidth]{images/dashboard1.jpg}
  \caption{Biểu đồ Radar và Bảng xếp hạng các cấu hình RAG}
  \label{fig:benchmark-radar}
\end{figure}

\subsection{Bảng So sánh Kết quả 4 Cấu hình RAG}

\begin{table}[H]
  \centering
  \caption{Kết quả Benchmark so sánh 4 cấu hình RAG (LLM cố định: Gemini 2.5 Flash)}
  \label{tab:benchmark-results}
  \renewcommand{\arraystretch}{1.6}
  \begin{tabularx}{\textwidth}{|l|c|c|c|c|}
    \hline
    \rowcolor{blue!15}
    \textbf{Cấu hình} & \textbf{Accuracy} & \textbf{Faithfulness} & \textbf{Relevancy} & \textbf{Latency (s)} \\
    \hline
    Recursive + Gemini Embed & \textbf{0.89} & \textbf{0.94} & \textbf{0.88} & 4.1 \\
    \hline
    Recursive + BAAI/bge-m3 & 0.85 & 0.92 & 0.86 & \textbf{3.2} \\
    \hline
    Fixed-size + Gemini Embed & 0.81 & 0.91 & 0.72 & 4.3 \\
    \hline
    Fixed-size + BAAI/bge-m3 & 0.78 & 0.90 & 0.70 & 3.4 \\
    \hline
  \end{tabularx}
\end{table}

\subsection{Phân tích Kết quả}

\subsubsection{Phân tích Chiến lược Phân mảnh (Chunking)}

Kết quả Benchmark cho thấy chiến lược \textbf{Recursive Chunking} áp đảo hoàn toàn Fixed-size Chunking về chỉ số Relevancy (đạt 0.88 so với 0.72 của Fixed-size). Nguyên nhân là:

\begin{itemize}[itemsep=3pt]
  \item \textbf{Fixed-size Chunking} cắt văn bản khi đạt giới hạn ký tự, bất kể cấu trúc ngữ pháp. Điều này dẫn đến các chunk bị ``đứt câu'' giữa chừng, làm mất ngữ cảnh quan trọng.
  \item \textbf{Recursive Chunking} ưu tiên cắt tại các điểm tự nhiên (dấu chấm câu, ngắt đoạn, tiêu đề). Các chunk giữ được trọn vẹn ý nghĩa ngữ nghĩa, giúp LLM có nền tảng tốt hơn để tổng hợp câu trả lời.
\end{itemize}

\subsubsection{Phân tích Mô hình Embedding}

\begin{itemize}[itemsep=3pt]
  \item \textbf{Google Gemini Embedding} cho Accuracy cao hơn $\approx 4\%$ nhờ khả năng hiểu ngôn ngữ tiếng Việt tự nhiên sâu sắc hơn (được huấn luyện trên tập dữ liệu đa ngôn ngữ khổng lồ).
  \item \textbf{BAAI/bge-m3} (chạy cục bộ) có Latency thấp hơn đáng kể (3.2s so với 4.1s), do không phải gọi API qua mạng, phù hợp cho môi trường không cần API key.
  \item Cả hai mô hình đều đạt Faithfulness xuất sắc ($> 0.90$), chứng tỏ ChromaDB truy xuất tốt các đoạn tài liệu có chứa thông tin cần thiết.
\end{itemize}

\subsubsection{Khuyến nghị Cấu hình Tối ưu}

\begin{tcolorbox}[
  enhanced, colback=green!5!white, colframe=green!60!black,
  fonttitle=\bfseries, title={Kết luận Benchmark -- Cấu hình Tối ưu}
]
  Cấu hình \textbf{Recursive Chunking + Google Gemini Embedding + Gemini 2.5 Flash LLM} cho hiệu năng tốt nhất về Accuracy (0.89) và Faithfulness (0.94), với Latency 4.1 giây -- nằm trong giới hạn chấp nhận được ($< 5$ giây theo NFR-01). \\[4pt]
  Nếu ưu tiên tốc độ phản hồi, \textbf{Recursive + BAAI/bge-m3} là lựa chọn thay thế với hiệu năng gần tương đương nhưng Latency thấp hơn 22\%.
\end{tcolorbox}

%------------------------------------------------------
\section{Tóm tắt Kết quả Kiểm thử}
%------------------------------------------------------

\begin{table}[H]
  \centering
  \caption{Tóm tắt kết quả kiểm thử toàn bộ hệ thống}
  \label{tab:test-summary}
  \renewcommand{\arraystretch}{1.5}
  \begin{tabularx}{\textwidth}{|l|c|c|c|}
    \hline
    \rowcolor{blue!10}
    \textbf{Hạng mục kiểm thử} & \textbf{Tổng TC} & \textbf{Pass} & \textbf{Tỷ lệ} \\
    \hline
    Kiểm thử chức năng & 20 & 20 & \textbf{100\%} \\
    \hline
    Kiểm thử phân quyền & 4 & 4 & \textbf{100\%} \\
    \hline
    Kiểm thử xử lý lỗi & 6 & 6 & \textbf{100\%} \\
    \hline
    Kiểm thử hiệu năng & 4 & 4 & \textbf{100\%} \\
    \hline
    RAGAS Benchmark (Faithfulness $\geq$ 0.85) & 4 configs & 4 & \textbf{100\%} \\
    \hline
    \rowcolor{green!10}
    \textbf{Tổng cộng} & \textbf{38} & \textbf{38} & \textbf{100\%} \\
    \hline
  \end{tabularx}
\end{table}

\subsection{Đánh giá Tổng kết}

Thông qua quá trình kiểm thử toàn diện, nhóm đã thu thập được các kết luận sau:

\begin{itemize}[itemsep=4pt]
  \item \textbf{Hệ thống đáp ứng 100\% Test Cases:} Toàn bộ 38 test cases đều Pass, khẳng định hệ thống hoạt động đúng theo đặc tả SWR.
  \item \textbf{Chất lượng AI vượt ngưỡng yêu cầu:} Cấu hình tốt nhất đạt Faithfulness = 0.94 $> 0.85$ (NFR-05) và Relevancy = 0.88 $> 0.80$ (NFR-06).
  \item \textbf{Hiệu năng đạt yêu cầu:} Latency trung bình 3.2 -- 4.3 giây, nằm trong giới hạn $< 5$ giây (NFR-01).
  \item \textbf{Hệ thống gần như loại bỏ hoàn toàn Hallucination:} Faithfulness $> 0.90$ chứng tỏ AI chỉ trả lời dựa trên tài liệu được cung cấp, không tự bịa đặt thông tin.
  \item \textbf{Benchmark là công cụ nghiên cứu giá trị:} Module Benchmark cung cấp số liệu định lượng khách quan, thay thế hoàn toàn việc đánh giá cảm quan.
\end{itemize}

\section{Kết luận Chương}

Chương này đã trình bày đầy đủ tài liệu kiểm thử (SWT) bao gồm: phạm vi kiểm thử (In-scope/Out-of-scope), chiến lược và mức độ kiểm thử, 20 Test Cases chức năng chi tiết, cấu hình môi trường Benchmark RAGAS, kết quả so sánh 4 cấu hình RAG với biểu đồ Radar và phân tích chuyên sâu, cùng bảng tóm tắt kết quả 100\% Pass. Kết quả kiểm thử khẳng định hệ thống Chatbot RAG đã đáp ứng đầy đủ và vượt trội các yêu cầu chức năng và phi chức năng được đặc tả trong SWR.
